\section*{Bab \ref{ch:g11:vect} --- Vektor dan Garis pada Bidang}

\begin{solution}{exo:g11:vect:1}
$\vect{AB}\,(3, 4)$, $\vect{AC}\,(-3, 2)$,
$\vect{AB} + \vect{AC}\,(0, 6)$,
$2\vect{AB} - 3\vect{AC}\,(6 + 9,\ 8 - 6) = (15, 2)$. \omterm{prop:g10:coordgeom:midpoint}{Titik tengah}
$[BC]$ adalah $\left(\frac{5 + (-1)}{2}, \frac{3+1}{2}\right) = (2, 2)$.
\end{solution}

\begin{solution}{exo:g11:vect:2}
$\det = 4 \times 3 - (-2) \times (-6) = 12 - 12 = 0$: \omterm{def:g11:vect:collinear}{kolinear} (memang
$\vec v = -\frac12 \vec u$).

$\det = 2 \times 7 - 3 \times 5 = -1 \neq 0$: tidak \omterm{def:g11:vect:collinear}{kolinear}.

$\det = 0 \times (-8) - 0 \times 3 = 0$: \omterm{def:g11:vect:collinear}{kolinear} (keduanya tegak).
\end{solution}

\begin{solution}{exo:g11:vect:3}
Melalui $A(2,1)$, arahnya $(3, -1)$: memadankan $(-b, a) = (3, -1)$ memberi
$a = -1$, $b = -3$, sehingga $-x - 3y + c = 0$; memasukkan $A$: $-2 - 3 + c =
0$, jadi $c = 5$. Setelah dikalikan $-1$: $x + 3y - 5 = 0$.

Melalui $B(0,4)$ dan $C(2,0)$: $\vect{BC}\,(2, -4)$, sehingga $a = -4$,
$b = -2$: $-4x - 2y + c = 0$; memasukkan $B$: $-8 + c = 0$, jadi $c = 8$.
Setelah disederhanakan dengan $-2$: $2x + y - 4 = 0$. Periksa dengan $C$: $4 + 0 - 4 = 0$.
\end{solution}

\begin{solution}{exo:g11:vect:4}
Sebuah \omterm{def:g11:vect:direction}{vektor arahnya} adalah $(-b, a) = (2, 3)$. Titik potongnya: $x = 0$
memberi $y = 3$, yaitu titik $(0, 3)$; $y = 0$ memberi $x = -2$, yaitu titik
$(-2, 0)$. Untuk $P(2,6)$: $3 \times 2 - 2 \times 6 + 6 = 0$, jadi $P \in d$.
Untuk $Q(1,4)$: $3 - 8 + 6 = 1 \neq 0$, jadi $Q \notin d$.
\end{solution}

\begin{solution}{exo:g11:vect:5}
$\vect{AB}\,(3, 6)$ dan $\vect{AC}\,(-3, -6)$:
$\det = 3 \times (-6) - (-3) \times 6 = 0$, jadi $A$, $B$, $C$ segaris.

$\vect{DE}\,(2, 3)$ dan $\vect{DF}\,(5, 8)$:
$\det = 2 \times 8 - 5 \times 3 = 1 \neq 0$: jadi $D$, $E$, $F$ tidak
segaris.
\end{solution}

\begin{solution}{exo:g11:vect:6}
Pasangan pertama: menjumlahkan kedua \omterm{def:g10:algebra:equation}{persamaannya} memberi $(2x + y - 7) + (x - y + 1) =
3x - 6 = 0$, jadi $x = 2$, lalu $y = x + 1 = 3$: keduanya bertemu di $(2, 3)$
(periksa: $4 + 3 - 7 = 0$).

Pasangan kedua: $ab' - a'b = 1 \times 4 - (-2) \times (-2) = 0$, sehingga
kedua garisnya sejajar; mengalikan $e$ dengan $-2$ memberi $-2x + 4y - 6 = 0$,
yang berbeda dari $e'$ ($-6 \neq 1$): jadi sejajar secara tegas, tanpa
titik potong.
\end{solution}

\begin{solution}{exo:g11:vect:7}
Untuk $d: 2x - y + 3 = 0$: arahnya $(-b, a) = (1, 2)$ dan titik
$(0, 3)$ terletak pada $d$, sehingga $d = \{(t,\ 3 + 2t) : t \in \R\}$.

Untuk garis parametrik melalui $(1, -3)$ dengan arah $(2, 1)$:
$a = 1$, $b = -2$, sehingga $x - 2y + c = 0$; memasukkan $(1, -3)$:
$1 + 6 + c = 0$, jadi $c = -7$. \omterm{def:g10:algebra:equation}{Persamaannya}: $x - 2y - 7 = 0$ (periksa $t = 1$,
titik $(3, -2)$: $3 + 4 - 7 = 0$).
\end{solution}

\begin{solution}{exo:g11:vect:8}
Garis yang sejajar mempunyai $(a, b)$ yang sama: $4x - 3y + c = 0$; melalui
$P(1, -2)$: $4 + 6 + c = 0$, jadi $c = -10$ dan garisnya adalah
$4x - 3y - 10 = 0$.

Kesejajaran $mx + 2y - 5 = 0$ dengan $d$ menuntut
$ab' - a'b = 4 \times 2 - m \times (-3) = 8 + 3m = 0$, jadi
$m = -\frac83$.
\end{solution}

\begin{solution}{exo:g11:vect:9}
\omterm{prop:g10:coordgeom:midpoint}{Titik tengah} $[AB]$ adalah $M\!\left(\frac{-1+3}{2}, \frac{0+2}{2}\right)
= (1, 1)$. Garis beratnya adalah garis $(CM)$ dengan $C(1, -4)$: kedua titiknya
berabsis $1$, sehingga garis beratnya adalah garis tegak
\[
x - 1 = 0 .
\]
\end{solution}

\begin{solution}{exo:g11:vect:10}
$\det(\vec u, \vec v) = m(m - 1) - 3 \times 2 = m^2 - m - 6$. Nilainya lenyap
ketika $m^2 - m - 6 = 0$: $\Delta = 1 + 24 = 25$, \omterm{def:g11:quad:discriminant}{akarnya}
$m = \frac{1 \pm 5}{2}$, yaitu $m = 3$ atau $m = -2$.
\end{solution}

\begin{solution}{exo:g11:vect:11}
\emph{1.} $I = \left(\frac12, 0\right)$. Lalu
$\vect{DI} = \left(\frac12 - 0,\ 0 - 1\right) = \left(\frac12, -1\right)$
dan $\vect{DJ} = \frac23\vect{DI} = \left(\frac13, -\frac23\right)$, sehingga
\[
J = D + \vect{DJ} = \left(\frac13,\ \frac13\right).
\]

\emph{2.} $\vect{AJ}\left(\frac13, \frac13\right)$ dan
$\vect{AC}\,(1, 1)$:
$\det = \frac13 \times 1 - 1 \times \frac13 = 0$, jadi $A$, $J$, $C$
segaris --- bahkan $\vect{AJ} = \frac13\vect{AC}$: garis $(DI)$ memotong
diagonal $(AC)$ tepat pada sepertiga panjangnya dari $A$.
\end{solution}

\begin{solution}{pb:g11:vect:1}
\textbf{1.} $x = 1 + 3t$, $y = 2 - t$. Dari yang kedua,
$t = 2 - y$; setelah disubstitusikan: $x = 1 + 3(2 - y) = 7 - 3y$:
jadi \omterm{def:g10:algebra:equation}{persamaan} Kartesiusnya $x + 3y - 7 = 0$.

\textbf{2.} Untuk $ax + by + c = 0$, sebuah \omterm{def:g11:vect:direction}{vektor arahnya} adalah
$(-b, a)$: di sini $(5, 2)$. Sebuah titiknya: $y = 1$ memberi $x = 1$:
yaitu $(1, 1)$. Parametriknya: $x = 1 + 5t$, $y = 1 + 2t$.

\textbf{3.} Dari $x = 7 - 3y$ pada $2x - 5y + 3 = 0$:
$14 - 6y - 5y + 3 = 0$, sehingga $y = \frac{17}{11}$ dan
$x = 7 - \frac{51}{11} = \frac{26}{11}$: yaitu titik
$\left(\frac{26}{11}, \frac{17}{11}\right)$.

\textbf{4.} Arahnya $(5, 2)$ dan $(3, -1)$:
$\det = 5 \times (-1) - 2 \times 3 = -11 \neq 0$: jadi berpotongan,
seperti yang sudah ditemukan. Untuk pasangan kedua: arahnya $(5, 2)$ dan
$(10, 4)$ mempunyai $\det = 0$ (sejajar), dan melipatduakan
$2x - 5y + 3 = 0$ memberi $-4x + 10y - 6 = 0 \neq -4x + 10y + 1
= 0$: jadi sejajar secara tegas.

\textbf{5.} $\vect{AB}(3, 2)$, $\vect{AC}(9, 6)$:
$\det = 18 - 18 = 0$: segaris.

\textbf{6.} $P$: dari $x = 2t$, $y = 3 + t$ diperoleh $y = 3 +
\frac x2$. $Q$: dari $x = 10 - t$ diperoleh $t = 10 - x$, sehingga
$y = 2(10 - x) - 1 = 19 - 2x$. Perpotongannya:
$3 + \frac x2 = 19 - 2x$ memberi $x = 6.4$, $y = 6.2$: jadi kedua lintasan
itu berpotongan di $(6.4,\ 6.2)$.

\textbf{7.} $P$ melewatinya ketika $2t = 6.4$: yaitu $t = 3.2$ jam. $Q$
ketika $10 - t = 6.4$: yaitu $t = 3.6$ jam. Terpaut dua puluh empat menit:
jadi tidak bertabrakan. Peringatannya: sebuah \omterm{def:g10:functions:function}{peta} menunjukkan
\emph{lintasan}; hanya jam parametriknya yang mengatakan apakah dua gerak
menempati titik yang sama pada saat yang sama.

\textbf{8.} $D^2(t) = (3t - 10)^2 + (4 - t)^2 =
10t^2 - 68t + 116$: titik puncaknya di $t = \frac{68}{20} = 3.4$ jam,
tempat $D^2 = 0.4$: jadi pendekatan terdekatnya
$D = \sqrt{0.4} \approx 0.63$ mil pada $t = 3.4$.

\textbf{9.} Dilanggar: $0.63 < 1$. Pertidaksamaan $D^2(t) < 1$ berbunyi
$10t^2 - 68t + 115 < 0$: $\Delta = 4624 - 4600 = 24$, \omterm{def:g11:quad:discriminant}{akarnya}
$\frac{68 \pm \sqrt{24}}{20} = 3.4 \pm 0.245$: jadi kedua kapal itu
terlalu dekat sejak $t \approx 3.16$ sampai $t \approx 3.64$ --- sekitar
$29$ menit pelanggaran. Salah satunya harus mengubah haluan.

\textbf{10.} Setiap \omterm{def:g10:coordgeom:system}{koordinat} setiap kapalnya merupakan \omterm{def:g11:func:function}{fungsi} \emph{afin}
dari $t$, sehingga setiap selisih \omterm{def:g10:coordgeom:system}{koordinatnya} pun afin,
dan $D^2$ --- jumlah dua kuadrat \omterm{def:g10:reffunc:affine}{fungsi afin} --- merupakan
bentuk kuadrat (dengan koefisien utama taknegatif). Karena itu pendekatan
terdekatnya selalu terbaca dari sebuah titik puncak.

\textbf{11.} Menemui $P$ pada $t = 3$, yaitu di $(6, 6)$:
$(3a, 3b) = (6, 6)$ memberi $(a, b) = (2, 2)$, dengan laju
$\sqrt{8} \approx 2.8$ knot.

\textbf{12.} $D(0, 1)$, $I\left(\frac12, 0\right)$: arahnya
$\left(\frac12, -1\right)$, yaitu $(1, -2)$: \omterm{def:g10:algebra:equation}{persamaannya}
$2x + y = 1$.

\textbf{13.} $(AC)$: $y = x$. Lalu $2x + x = 1$:
$x = \frac13$: yaitu titik $\left(\frac13, \frac13\right)$, pada
sepertiga jalan dari $A$ ke $C$.

\textbf{14.} $B(1, 0)$, $K\left(\frac12, 1\right)$: arahnya
$\left(-\frac12, 1\right)$, \omterm{def:g10:algebra:equation}{persamaannya} $2x + y = 2$. Dengan
$y = x$: $x = \frac23$: yaitu titik
$\left(\frac23, \frac23\right)$. Kedua garis ceva itu memotong $(AC)$ di
kedua titik pembagi tiganya: teoremanya terbukti, dengan dua kali tiga baris
hitungan.

\textbf{15.} Setiap jajargenjang dapat dibawa ke persegi satuan
dengan memilih $A$ sebagai titik asal serta $\vect{AB}$ dan $\vect{AD}$ sebagai
satuan kedua sumbunya. Pilihan itu mempertahankan segala yang dibicarakan
teoremanya --- \omterm{prop:g10:coordgeom:midpoint}{titik tengah}, kesegarisan, dan perbandingan sepanjang sebuah
garis (semuanya dapat dinyatakan dengan \omterm{def:g11:vect:vector}{vektor} dan skalar) --- sehingga
membuktikan pernyataannya pada persegi itu sudah membuktikannya untuk setiap
jajargenjang. (Pilihan itu \emph{tidak} membawa panjang atau sudut: keduanya
diubah oleh kerangkanya.)

\textbf{16.} $E\left(1, \frac12\right)$: garis $(AE)$:
$y = \frac x2$. Diagonal $(DB)$: dari $(0,1)$ ke $(1,0)$:
$y = 1 - x$. Perpotongannya: $\frac x2 = 1 - x$: $x = \frac23$:
yaitu titik $\left(\frac23, \frac13\right)$ --- sekali lagi sebuah
titik pembagi tiga, kini pada diagonal yang lain. Garis ceva lewat \omterm{prop:g10:coordgeom:midpoint}{titik
tengah} memang gemar sepertiga.

\textbf{17.} Dua yang pertama: $x + y = 4$ dan $2x - y = 2$ memberi
$3x = 6$: yaitu $(2, 2)$. Yang ketiga: $2 - 4 = -2$: terpenuhi. Jadi
ketiganya berpotongan di satu titik, $(2, 2)$.

\textbf{18.} $mx - y + 2 - 3m = m(x - 3) + (2 - y) = 0$: agar
\omterm{def:g10:algebra:equation}{persamaannya} berlaku untuk \emph{setiap} $m$, ambillah $x = 3$,
$y = 2$ --- dan memang setiap $d_m$ memuat $(3, 2)$. Jadi sebuah berkas
garis melalui satu titik, satu garis untuk setiap \omterm{def:g10:reffunc:affine}{gradien} $m$.

\textbf{19.} Sejajar dengan $y = 2x + 7$: \omterm{def:g10:reffunc:affine}{gradiennya} $m = 2$, yaitu garis
$2x - y - 4 = 0$. Melalui titik asal: $2 - 3m = 0$:
$m = \frac23$.

\textbf{20.} Kartesius: masukkan sebuah titik, keanggotaannya terjawab.
Tereduksi: \omterm{def:g10:reffunc:affine}{gradien} dan titik potong sumbunya terlihat, \omterm{def:g10:functions:graph}{grafiknya} seketika.
Parametrik: jamnya sudah terpasang --- tabrakan, pendekatan
terdekat, pencegatan. \omterm{def:g11:vect:det}{Determinan}: satu perkalian
menyilang, dan kesejajarannya diputuskan. Kerangka yang dipilih: sebuah
jajargenjang menjadi persegi satuan, dan sebuah teorema pembagian tiga yang
klasik menjadi tiga perpotongan garis.

\end{solution}
